% s-029-p029-m002; original PDF 29, 30, 31, 32
\Aphorism{s-029-p029-m002}{20}
Der geistige Fortschritt besteht weit seltener in der Auffindung neuer Gedanken und Formeln, als in der sinnvollen Deutung und Beseelung der alten, bisher ohne Bewusstsein ausgesprochenen und überlieferten. Goethes Zweifel „Wer kann was Dummes, wer was Kluges denken, das nicht die Vorwelt schon gedacht?“ ist gewiss begründet: nur bleibt hinzuzufügen, dass die Menschheit nichts lieber zu thun pflegt, als Dinge zu denken, Fragen zu stellen, Begriffe zu formuliren, auf die sie nach ihrem augenblicklichen Zustande durchaus kein Recht hat — dieses Recht nachträglich hinzuzuerwerben ist Sache der Nachwelt. Gedacht und geredet wird zu allen Zeiten dasselbe; aber der reale Hintergrund, vor dem sich diese dialectische Komödie abspielt, die Vorgeschichte sinnlicher und leidenschaftlicher Erfahrungen, auf die hin jenes Denken und Reden zu Recht besteht und ohne die es nur vorlaute Kinderschulweisheit bleiben würde — diese rechtfertigende Grundlage und Vorarbeit wird von Jahrhundert zu Jahrhundert eine andere, und wie es scheint, eine festere, tiefere, umfangreichere, werthvollere. Danach ist es gründlich verfehlt (abgesehen davon, dass es hochmüthige und doch wohlfeile Weisheit ist), wenn überkluge Historiomanen irgend einer neuen Erscheinung damit begegnen wollen, dass sie aus ihrem Wissensvorrath etwas Ähnliches hervorsuchen und nun von Wiederaufnahme eines „überwundenen Standpunktes“ reden. Was \Foreign{ad oculos} ähnlich aussieht, ist \Foreign{ad rem} noch nicht dasselbe, und wenn zwei Menschen aus getrennten Culturen das Gleiche sagen , so meinen sie doch sicher nicht das Gleiche; man muss sich nicht von Buchstaben und Augenschein übertölpeln lassen. Nehmen wir ein beliebiges Beispiel. Schopenhauer kann es nicht lassen, die moderne \Sperr{materialistische} Naturauffassung immer wieder mit dem Hinweis auf unglückliche Vorgänger, wie Demokrit und die Atomistiker, abzufertigen. Nun braucht man noch lange kein Freund der materialistischen Oberflächlichkeit und Alleswisserei zu sein, aber soviel einzugestehen gebietet doch die Gerechtigkeit, dass unser moderner Materialismus eine ganz andere Vergangenheit hinter sich und also ein ganz anderes Existenzrecht für sich hat als die speculative Systembaukunst der griechischen Philosophen. Demokrits Atome haben mit der Wissenschaft soviel zu thun wie der Nous des Anaxagoras oder die Urstoffe der Milesier, nämlich beinahe nichts; es sind mythopoetische Annahmen, die für die Erklärung und Beschreibung von Naturvorgängen höchstens eine grosse Voreiligkeit bedeuten würden. Unser heutiger Materialismus hingegen kommt immerhin von einigen Jahrhunderten Arbeit, Empirie, Analysis her, und seine Darstellung ist für die untersten Arten des Geschehens, für die rein physicalischen Erscheinungen, im Ganzen befriedigend. Dass er diese Darstellung auch auf die höheren Arten, auf physiologische, sociale, psychische Erscheinungen ausgedehnt wissen will, ist wiederum \Sperr{seine} Voreiligkeit, aber keine schlimmere, als deren sich alle angeblichen „Erklärungen“ geistiger Phänomene schuldig gemacht haben. Für einen billig Denkenden ist der Materialist, der diese Phänomene zwar unter seinen Problemkreis rechnet, aber bereitwillig ihre ausserordentlich hohe Complicationsstufe anerkennt, immer noch ein Muster an Takt und ehrerbietigem Naturempfinden gegenüber dem Heros der „Geisteswissenschaften“, der mit „frecher Stirne, kühner Brust“ den Menschen in Sinnlichkeit, Verstand und Vernunft eintheilt und ohne Herzklopfen den Schmetterling „Seele“ auf eine Begriffsnadel spiesst. — Die retrospectiven Analogiensucher mögen sich also vorsehen, wenn sie auf Grund von Gesichts- und Namensähnlichkeiten Zwei für blutsverwandt halten, die durch Raum und Zeit geschieden sind; und dass ein „überwundener Standpunkt“ bei erneuertem Auftreten immer wieder überwindbar sei, dieser Aberglaube sollte zuallererst als überwundener Standpunkt gelten. Damit, dass seine ersten Voreiligkeiten und Entwicklungsansätze ein unglückliches Ende nehmen, ist ein Ideal, ein Gedanke noch durchaus nicht beseitigt; im Gegentheil, ein lebensfähiger Culturorganismus gewinnt um so mehr an Kraft und Keimtrieb, je schonungloser alle verfrühten und schwächlichen Schösslinge ausgerottet wurden. Das Dumme, das die Vorwelt gedacht, kann eines Tages noch klug werden, wenn sich das nöthige Erfahrungssubstrat darunter gelegt hat; der Weltgeschichte, die als Weib immer bei boshafter Laune ist, macht es Vergnügen, jedem nothwendigen und dauerhaften Typus eine embryonale Halbheit voranlaufen zu lassen, die zur Unzeit ihr unreifes Dasein durchsetzen will und darüber zu Grunde geht.
\par

